% The Trial: every number comes from build/the-trial/numbers.tex (\lab{...})
\input{../build/the-trial/numbers.tex}
\graphicspath{{../build/the-trial/}}
\newcommand\strip[2]{% #1 measure key, #2 title
  \begin{tikzpicture}
  \begin{axis}[width=40mm, height=21mm, scale only axis=false, ymin=0, axis y line=none,
      xmin=\lab{bmin-#1}, xmax=\lab{bmax-#1}, enlarge x limits=0.06, xtick=\empty,
      title={#2}, title style={font=\figfont\scriptsize, yshift=-1.5mm},
      clip=false]
    \addplot[const plot mark mid, black, line width=0.4pt] table {../build/the-trial/baseline-#1.dat};
    \draw[line width=1.1pt] (axis cs:\lab{raw-#1},0) -- (axis cs:\lab{raw-#1},\pgfkeysvalueof{/pgfplots/ymax});
    \node[anchor=north, font=\figfont\scriptsize] at (axis cs:\lab{raw-#1},0) {$z=\lab{z-#1}$};
  \end{axis}
  \end{tikzpicture}}

\chapter[The Trial]{The Trial}
\label{ch:trial}

\begin{center}\small
\textsc{Franz Kafka}, \emph{Der Proce\ss}, 1925\\
in the English translation by \lab{translator} (Project Gutenberg \#\lab{gid})
\end{center}

\bigskip
\noindent The specimen has \lab{sentences} sentences and \lab{words} words in
\lab{chapters} chapters. \lab{m-share-left-pct}\,\% of its sentences turn left
(positive), \lab{m-share-right-pct}\,\% right (negative), and
\lab{m-share-straight-pct}\,\% go straight. Overall the book turns right
\lab{m-handedness-abs} full turns more than it turns left. A sentence averages
\lab{m-stride_mean} words, with a coefficient of variation of
\lab{m-stride_cv}.

Those numbers describe what the book is made of. The table describes its order:
each measure set against the same sentences shuffled \lab{shuffles} times.

\begin{table}[h]
\centering\figfont\small
\begin{tabular}{@{}lrrr@{}}
\toprule
 & \emph{The Trial} & shuffled & $z$\\
\midrule
Recoil (same-way turn pairs) & \lab{m-recoil} & \lab{b-recoil} $\pm$ \lab{sd-recoil} & \textbf{\lab{z-recoil}}\\
Vacillation (opposite pairs) & \lab{m-vacillation} & \lab{b-vacillation} $\pm$ \lab{sd-vacillation} & \textbf{\lab{z-vacillation}}\\
Turn entropy $H_2$ (bits) & \lab{m-entropy_h2} & \lab{b-entropy_h2} $\pm$ \lab{sd-entropy_h2} & \textbf{\lab{z-entropy_h2}}\\
Mean corridor (sentences) & \lab{m-corridor_mean} & \lab{b-corridor_mean} $\pm$ \lab{sd-corridor_mean} & \lab{z-corridor_mean}\\
Longest corridor & \lab{m-corridor_max} & \lab{b-corridor_max} $\pm$ \lab{sd-corridor_max} & \lab{z-corridor_max}\\
Drift & \lab{m-drift} & \lab{b-drift} $\pm$ \lab{sd-drift} & \lab{z-drift}\\
Gyration & \lab{m-gyration_final} & \lab{b-gyration_final} $\pm$ \lab{sd-gyration_final} & \lab{z-gyration_final}\\
Sprawl & \lab{m-sprawl} & \lab{b-sprawl} $\pm$ \lab{sd-sprawl} & \lab{z-sprawl}\\
Texture (dimension) & \lab{m-texture} & \lab{b-texture} $\pm$ \lab{sd-texture} & \lab{z-texture}\\
Crossings per sentence & \lab{m-crossing_rate} & \lab{b-crossing_rate} $\pm$ \lab{sd-crossing_rate} & \lab{z-crossing_rate}\\
Median crossing gap & \lab{m-crossing_gap_median} & \lab{b-crossing_gap_median} $\pm$ \lab{sd-crossing_gap_median} & \lab{z-crossing_gap_median}\\
Palimpsest (words) & \lab{m-palimpsest_words} & \lab{b-palimpsest_words} $\pm$ \lab{sd-palimpsest_words} & \lab{z-palimpsest_words}\\
\bottomrule
\end{tabular}
\caption*{Vital signs. Bold: significant ($|z|\ge3$, see p.~\pageref{ch:method}).}
\end{table}

\clearpage
\section{Reading the drawing}

The drawing opposite is the whole book. It begins at the black square and ends
at the open one. Open circles on the path mark where chapters begin. Every
label points to a structure found by the procedures of the method chapter. None
was placed by hand.

\begin{finding}{Knots I and II} The two densest tangles lie side by side near
the top of the drawing. Knot~I has \lab{knotA-crossings} crossings among
\lab{knotA-nsent} sentences, from \lab{knotA-first} to \lab{knotA-last}, in
chapters \lab{knotA-chapters}. They are the painter Titorelli's account of
acquittal (chapter~7), the businessman Block's account of his own case
(chapter~8, \lab{knotA-topn} of the sentences) and K.'s arrival in the cathedral
(chapter~9). Three scenes that are hundreds of sentences apart in the book are
drawn over one another in one place.
\end{finding}

\begin{finding}{Hub} The most-crossed sentence is \lab{hub-ref}, crossed
\lab{hub-n} times. It is Block's: ``\lab{hub-quote}'' Eleven of the thirteen
sentences that cross it come from the same conversation with Block in
chapter~8, and two from the cathedral. The sentence that describes a case going
round in circles is the one the path keeps returning to.
\end{finding}

\begin{finding}{Coil} From \lab{coil-first} to \lab{coil-last} the path turns
\lab{coil-q} quarter turns to the \lab{coil-dir}, two and a half full circles
in twelve sentences. It is the lawyer Huld's report of progress: ``the results
are very encouraging'', ``nothing has been lost yet''. The positive words make
it a left-hand spiral. The narration says ``there was always some progress, but
he could never be told what sort of progress it was''. The drawing records the
surface of the talk, which literally goes round in circles.
\end{finding}

\begin{finding}{Corridor} The longest run of neutral sentences
(\lab{cor-n}~sentences, \lab{cor-words} words, chapter~\lab{cor-ch}) is the
inventory of Titorelli's room: ``It was hardly long enough or broad enough to
make two steps.'' Description without feeling draws a straight line.
\end{finding}

\begin{finding}{Darkest passage} The thirty sentences with the lowest summed
score (\lab{dark-score}) begin at \lab{dark-first}, in chapter~1. It is not
the execution but Mrs.~Grubach's insinuations about Miss Bürstner, and K.'s
angry reply.
\end{finding}

\begin{finding}{Return} \lab{ret-i} (chapter~\lab{ret-chi}) and \lab{ret-j}
(chapter~\lab{ret-chj}) cross, \lab{ret-gap} sentences apart: ``\lab{ret-qi}''
and, on the way to the execution, ``\lab{ret-qj}''
\end{finding}

\begin{finding}{Palimpsest} The longest overwriting is \lab{pal-len} words of
\lab{pal-j} laid on \lab{pal-i}, both in chapter~\lab{pal-chi}. The drawing has
\lab{pal-n} such overlaps.
\end{finding}

\clearpage
\thispagestyle{empty}
\noindent\hspace*{0mm}\labelledfigure{maze.pdf}{../build/the-trial/maze-labels.tex}

\clearpage
\section{Against chance}

Each panel shows one measure: the outline is its distribution over
\lab{shuffles} shuffles of \emph{The Trial}'s own sentences, and the bar is
the book.

\medskip
\noindent\strip{recoil}{Recoil}\hfill\strip{vacillation}{Vacillation}\hfill\strip{entropy_h2}{Turn entropy $H_2$}\\[1mm]
\strip{corridor_mean}{Mean corridor}\hfill\strip{corridor_max}{Longest corridor}\hfill\strip{drift}{Drift}\\[1mm]
\strip{gyration_final}{Gyration}\hfill\strip{sprawl}{Sprawl}\hfill\strip{texture}{Texture}\\[1mm]
\strip{crossing_rate}{Crossings / sentence}\hfill\strip{crossing_gap_median}{Median crossing gap}\hfill\strip{palimpsest_words}{Palimpsest}

\medskip
Three measures leave the shuffled range: recoil, vacillation and $H_2$. All
three say the same thing. \emph{The Trial}'s feelings come in runs. A negative
sentence is followed by another negative one, and a positive by a positive,
more often than chance allows. A turn is then partly predictable from the turn
before. The table counts consecutive pairs (rows: first sentence, columns:
second; L~left, S~straight, R~right) against the number expected if order did
not matter:

\begin{center}\figfont\small
\begin{tabular}{@{}lrrr@{}}
\toprule
 & L & S & R\\
\midrule
L & \lab{tr-LL} / \lab{te-LL} & \lab{tr-LS} / \lab{te-LS} & \lab{tr-LR} / \lab{te-LR}\\
S & \lab{tr-SL} / \lab{te-SL} & \lab{tr-SS} / \lab{te-SS} & \lab{tr-SR} / \lab{te-SR}\\
R & \lab{tr-RL} / \lab{te-RL} & \lab{tr-RS} / \lab{te-RS} & \lab{tr-RR} / \lab{te-RR}\\
\bottomrule
\end{tabular}
\end{center}

\noindent Same-way pairs are over-represented (L\,L \lab{tq-LL}$\times$,
R\,R \lab{tq-RR}$\times$), and reversals of feeling are under-represented
(R\,L \lab{tq-RL}$\times$). In the drawing this shows as recoil: the path
doubles back on itself, which is why it is so dense. The large-scale shape,
however, is not distinguishable from chance. Drift, gyration, sprawl, texture
and the number of crossings all fall inside the shuffled range. Any ordering
of these sentences would have tangled about this much. What the order adds is
\emph{where} the tangle sits and which sentences it ties together.

\clearpage
\section{Knot I, enlarged}

\begin{figure}[h]
\centering
\includegraphics{zoom-A.pdf}
\caption{The square of $2\times\lab{knotA-radius}$ word units around Knot~I.
\lab{knotA-legend}. Open circles: crossings.}
\end{figure}

\noindent Knot~I is made of \lab{knotA-nsent} sentences, and its crossings join
sentences a median \lab{knotA-gap} sentences apart. The tangle is not one
passage folding on itself. Three scenes meet in one place: the painter's
account of acquittal, Block's years-long case and K.'s arrival in the empty
cathedral. Knot~II, beside it (\lab{knotB-crossings}
crossings, chapters \lab{knotB-chapters}), is made of the same material in
different proportions (\lab{knotB-legend}).

\clearpage
\section{Time}

\begin{center}
\begin{tikzpicture}[x=110mm,y=1mm, font=\figfont\scriptsize]
  % knot I sentences over narrative time
  \begin{scope}[yshift=10mm]
  \pgfplotstableread{../build/the-trial/knot-time.dat}\kt
  \pgfplotstableforeachcolumnelement{y}\of\kt\as\cell{%
    \pgfplotstablegetelem{\pgfplotstablerow}{x}\of\kt
    \fill (\pgfplotsretval-0.0125,0) rectangle ++(0.025,\cell*0.5);}
  \node[anchor=west] at (1.01,1) {Knot I};
  \end{scope}
  \draw[line width=0.4pt] (0,0) -- (1,0);
  \edef\chl{0,\lab{tl-chapters},1}
  \edef\chs{0,\lab{tl-chapters}}
  \foreach \x in \chl {\draw[line width=0.4pt] (\x,-1.2) -- (\x,1.2);}
  \foreach \x [count=\c] in \chs {\node[anchor=north west, inner sep=1pt] at (\x,-1.2) {\c};}
  \node[anchor=west] at (1.01,0) {chapters};
  % events
  \edef\evs{\lab{tl-coil}/coil,\lab{tl-dark}/darkest,\lab{tl-cor}/corridor,\lab{tl-hub}/hub}
  \foreach \p/\t in \evs
     {\draw (\p,-6) -- (\p,-9); \node[anchor=north, rotate=0] at (\p,-9) {\t};}
  \draw (\lab{tl-ret-i},-15) to[bend right=10] node[midway,below] {return} (\lab{tl-ret-j},-15);
  \fill (\lab{tl-ret-i},-15) circle (0.5mm); \fill (\lab{tl-ret-j},-15) circle (0.5mm);
\end{tikzpicture}
\end{center}

\noindent On the time line, every structure in the drawing is placed at its
position in the book. Knot~I's sentences come in two bursts, the painter and
Block, then a last few in the cathedral. The coil, the corridor and the hub
all fall in the second half, in chapters 7 and 8, where K. is passed from one
adviser to the next. The darkest passage comes much earlier, in chapter~1.

\begin{center}
\begin{tikzpicture}
\begin{axis}[width=105mm, height=48mm, xmin=0, xmax=1, ymin=0, xlabel={narrative time}, ylabel={radius of gyration (words)},
  label style={font=\figfont\scriptsize}, legend style={draw=none, font=\figfont\scriptsize, at={(0.02,0.98)}, anchor=north west}]
  \addplot[name path=hi, draw=none, forget plot] table[x=t,y=hi] {../build/the-trial/gyration.dat};
  \addplot[name path=lo, draw=none, forget plot] table[x=t,y=lo] {../build/the-trial/gyration.dat};
  \addplot[pattern=north east lines, draw=none] fill between[of=hi and lo];
  \addlegendentry{shuffled, 5--95\,\%}
  \addplot[black, dashed, line width=0.4pt] table[x=t,y=mid] {../build/the-trial/gyration.dat};
  \addlegendentry{shuffled, median}
  \addplot[black, line width=1pt] table[x=t,y=rg] {../build/the-trial/gyration.dat};
  \addlegendentry{\emph{The Trial}}
\end{axis}
\end{tikzpicture}
\end{center}

\noindent How far the drawing has spread at each moment of reading: half its
final spread is reached after \lab{rg-half}\,\% of the book. From about the
middle of the book to nine tenths of the way through, \emph{The Trial} has
spread farther than 95\,\% of its shuffles (\lab{rg-above-pct}\,\% of
narrative time in all). In the last tenth it draws back in, and its final
spread is within the shuffled range ($z=\lab{z-gyration_final}$). The curve
compares a hundred moments at once, so it is descriptive, not a test. It shows
a book that strays farthest from its starting ground just before the end.

\clearpage
\section{The mineshaft}

\begin{figure}[!h]
\centering
\labelledfigure{mineshaft.pdf}{../build/the-trial/mineshaft-labels.tex}
\caption{\emph{The Trial} as a shaft. Each left turn tilts the path
\lab{shaft-delta}\textdegree{} up and each right turn tilts it
\lab{shaft-delta}\textdegree{} down. Seen from an angled side view, in
perspective; numbers mark chapter starts; vertical scale \lab{shaft-vscale}.}
\end{figure}

\noindent The shaft makes handedness visible. Because \emph{The Trial} turns
right more than left, it descends. Its tilt reaches \lab{shaft-min-pitch}\textdegree{}
and is still \lab{shaft-final-pitch}\textdegree{} at the last stroke. The deepest
point of the whole book is its last sentence: ``\lab{shaft-deepest-quote}''

\clearpage
\section{Findings}

\begin{enumerate}
\item \emph{The Trial}'s feelings persist. Recoil (\z{recoil}), vacillation
(\z{vacillation}) and conditional turn entropy (\z{entropy_h2}) all depart
from chance. Consecutive sentences agree in valence more often than the same
sentences in random order.
\item Its large-scale geometry is that of its material. Drift, gyration,
sprawl, texture and crossing rate are all within $|z|<3$. The drawing is as
tangled as any shuffle of these sentences.
\item Its densest tangles join chapters 7--9: the painter, Block and the
cathedral. The most-crossed sentence is Block's description of a case that
repeats itself. (Descriptive; knot positions are not tested against chance.)
\item The lawyer's reassurances form the drawing's tightest spiral, a positive
one. The word list reads the surface of a speech the novel presents as empty.
This is a limit of the method that the drawing makes visible.
\item In three dimensions the book descends, and it ends at its lowest point.
\end{enumerate}
